\documentclass[11pt]{article}
\usepackage[a4paper,margin=2.2cm]{geometry}
\usepackage{booktabs,array,longtable,xcolor,iftex}
\ifXeTeX\usepackage{fontspec}\setmainfont{texgyretermes}[Extension=.otf, UprightFont=*-regular, BoldFont=*-bold, ItalicFont=*-italic, BoldItalicFont=*-bolditalic]\fi
\setlength{\parindent}{0pt}\setlength{\parskip}{5pt}
\begin{document}
{\Large\bfseries Response to the review of the first version\par}\medskip
Stage-Wise Attribution of Geometric Error in Automatic Anatomical Surface Reconstruction From Computed Tomography\par
\textit{Sameer Morya and Sarthak Wage, K J Somaiya School of Engineering, Mumbai}\par\medskip
Thank you for the detailed review. We addressed each point with new experiments or analyses where the data allowed, and with tempered claims where they did not. The revision adds a held-out second organ (MSD Task07 Pancreas, 41 scans drawn at random) tested against six hypotheses written down before any pancreas scan was processed (\texttt{research/heldout\_hypotheses.md}); TotalSegmentator's full-resolution mode as a second engine configuration; a coarsening-model sensitivity analysis; re-scoring of every mesh with local and task-specific metrics; and a 12-section supplement (S1–S12). The paper remains six pages; the first version's Tier 1 table and its two resolution and topology figures moved to the supplement.\par
The held-out test changed some conclusions, and the paper now says so. The sampling relationship, the dominance of spacing for correct segmentations, the dominance of scans for the engine and the thick-slice penalty transferred to the pancreas (H1, H2, H3, H5 confirmed). The spleen's insensitivity to 3 mm grids did not (H4 not confirmed), and meshing on the engine's 3 mm grid improved closure but cost accuracy on the pancreas (H6 partly confirmed). That insensitivity is also specific to the fast engine mode: full-resolution mode is more accurate on native spleen scans but degrades by 0.42 mm on 3 mm input.\par
Re-scoring the spleen with the extended code reproduced all 16,332 result rows of the first version exactly, so every number carried over from the first version is unchanged.\par
\section*{Weaknesses}
\begin{longtable}{@{}>{\raggedright\arraybackslash}p{0.22\textwidth}>{\raggedright\arraybackslash}p{0.56\textwidth}>{\raggedright\arraybackslash}p{0.16\textwidth}@{}}
\toprule Review point & What we did & Where \\ \midrule\endhead
\textbf{W1} Limited external validity: one organ, one dataset, one annotator, one engine mode; thin, branching or pathological structures untested. & Added the pancreas, a thin, lobulated organ whose labels include tumours, as a held-out organ with six pre-specified hypotheses: four confirmed, one partly, one not (Table IV). Added TotalSegmentator full-resolution mode on both organs (S10). Still open and stated in the Limitations: one annotator per scan, both organs from one centre (MSKCC, through the MSD), no branching structure. & Sec. V, VI-C, VI-E, IX; Table IV; S10 \\ \midrule
\textbf{W2} Simulated coarsening does not reproduce kernel, noise or partial volume; 0.177$\Delta$ likely a lower bound and protocol-specific. & Repeated the coarse levels with a rectangular slice profile (a slab average) and with 20 HU noise added to the CT. The expert mask's error changed by at most 0.003 mm (slope 0.174 vs 0.177) and the engine's by at most 0.020 mm. The paper now limits the relationship to the tested protocol and calls it a probable lower bound; real thin and thick reconstructions of the same acquisition are not available in the MSD. & Sec. IV, VII, IX; S9 \\ \midrule
\textbf{W3} Overclaimed scope: “all stages”. & The contributions now read “crosses the tested reconstruction stages with two segmentation sources, an expert mask and one engine in two modes”, and no sentence in the abstract or elsewhere claims all stages. & Abstract; Sec. I \\ \midrule
\textbf{W4} $\eta^2$ presented too causally; case difficulty inferred, not measured. & $\eta^2$ is defined as descriptive in Sec. III and called descriptive wherever it is used. “Case difficulty” is replaced by “between-scan variation”. S11 measures scan properties: low organ-to-surround contrast goes with higher error on both organs ($\rho$ = --0.59, --0.49), but the properties explain little of the spleen's between-scan variance (leave-one-out R$^2$ 0.14) and more of the pancreas's (0.59), where tumour burden matters. & Sec. III, VII, VIII, IX; S11 \\ \midrule
\textbf{W5} The empirical law is only moderately predictive (R$^2$ $\approx$ 0.75; slice-only error 0.109 mm). & Added a 95\% prediction interval ($\pm$0.111 mm) and compared three refinements: a free intercept, linear axis terms and quadrature axis weights. None generalises better; the quadrature weights converge to equal values (0.031, 0.032), which supports the pooled $\Delta$. On the held-out pancreas the spleen fit has a mean absolute error of 0.030 mm and its interval covers 98\% of points. Slice-only coarsening remains the weakest prediction, and the paper says so. & Sec. VI-B; Fig. 2; S4 \\ \midrule
\textbf{W6} The slice5 subset has only 14 scans. & S5 gives per-scan effects (all 14 rose), the leave-one-scan-out range (0.175–0.195 mm for the engine) and the minimum detectable effect at 80\% power (0.063 mm against an observed 0.188 mm). The pancreas adds 35 resampled scans and replicates the effect (+0.143 mm, 34 of 35 scans). & Sec. VII; Table II; S5 \\ \midrule
\textbf{W7} Global metrics only; local defects and task tolerances not analysed. & Every mesh was re-scored with NSD at 0.5, 1 and 2 mm, the 99th percentile, the area of local defects (faces more than 2 mm from the reference) and their largest connected patch, volume error and two illustrative tolerances. Defects cover 8.9\% of the fast engine's spleen surface and 29\% of the pancreas surface; error maps and the largest missed or added region of each failure case are reported. & Sec. VI-D; S3, S8 \\ \midrule
\textbf{W8} Blur in voxels (production) vs millimetres (research) complicates comparisons. & Added a production-mm configuration (production smoothing and decimation, blur 0.5 mm) at every level. It changes production results by at most 0.065 mm on expert masks and 0.14 mm on engine masks; the grid comparison (H6) already holds blur in millimetres. & Sec. VI-D; S7 \\ \midrule
\textbf{W9} Phantom interactions may not generalise to anatomy. & S2 tests the smoothing-by-shape interaction on both organs. The reversal is shape-specific: no organ showed the capsule's doubling of error, though the pancreas is more sensitive to smoothing than the spleen. & S2 \\ \midrule
\textbf{W10} One field-of-view-cut scan caps the topology curves. & S6 repeats every closure rate without that scan: every undecimated mesh is then closed, and the effect of decimation is unchanged. & Sec. VI-D; S6 \\ \midrule
\textbf{W11} No held-out confirmation of the 3 mm working-grid finding. & Tested as H6 on the pancreas: the closure gain at 0.75 decimation held (+25 points [17, 32]), but accuracy did not (+0.085 mm [0.034, 0.143]), so H6 is partly confirmed and the recommendation now says to check the accuracy cost for thin organs. & Sec. VI-E, VIII; Table IV \\ \midrule
\textbf{W12} Production settings evaluated narrowly. & S12 recomputes every engine conclusion separately for each of the 49 configurations. Scan dominance and the thick-slice penalty hold in all of them for both organs and modes; the pancreas's 3 mm penalty (H4) is positive in every configuration, so it is not an artefact of one setting. & Sec. VII; S12 \\ \bottomrule
\end{longtable}
\section*{Suggested improvements}
\begin{longtable}{@{}>{\raggedright\arraybackslash}p{0.22\textwidth}>{\raggedright\arraybackslash}p{0.56\textwidth}>{\raggedright\arraybackslash}p{0.16\textwidth}@{}}
\toprule Review point & What we did & Where \\ \midrule\endhead
\textbf{I1} Temper claims; label $\eta^2$ as descriptive. & Done; see W3 and W4. & Abstract; Sec. I, III, VIII, IX \\ \midrule
\textbf{I2} Held-out or multi-organ validation. & Done with the pancreas and pre-specified hypotheses; see W1 and W11. & Sec. V, VI-E; Table IV \\ \midrule
\textbf{I3} Real coarse acquisitions or realistic noise, kernel and partial-volume simulation. & Partly: two simulation variants (slice profile, noise) bracket the effect (W2). Real paired thin and thick reconstructions are not in the dataset; listed below as not done. & S9 \\ \midrule
\textbf{I4} Multiple annotators or engines/modes. & Partly: a second engine mode was added (W1, I10). The MSD provides one annotation per scan, so multiple annotators were not possible. & S10 \\ \midrule
\textbf{I5} Refine the model; validate slice-only coarsening; prediction intervals. & Done; see W5. & Sec. VI-B; S4 \\ \midrule
\textbf{I6} Task-specific and local metrics; thin structures; tolerances. & Done; see W7. Thin structures: the Tier 1 capsule and the pancreas. Tolerances are illustrative, not clinically validated. & Sec. VI-D; S8 \\ \midrule
\textbf{I7} Clarify design and counts; which comparisons use production vs canonical. & S1 tabulates units, configurations, engine calls and meshes for every run (57,944 meshes) and states which analysis uses which settings. & Sec. V; S1 \\ \midrule
\textbf{I8} Strengthen the slice5 analysis. & Done; see W6. & S5 \\ \midrule
\textbf{I9} Matched physical-blur comparisons. & Done; see W8. & S7 \\ \midrule
\textbf{I10} More engine configurations, including full-resolution mode. & Added on both organs. On the spleen, full mode is more accurate at native resolution (ASSD 0.790 vs 0.901 mm, 39 of 41 scans better) but degrades on 3 mm input (+0.420 mm); on the pancreas it is worse than fast mode. Resolution insensitivity is therefore fast-mode-specific, while scan dominance holds for both modes. & Sec. VI-C; Table II; S10 \\ \midrule
\textbf{I11} Per-scan distributions and failure cases. & S3: a per-scan figure for every organ, source and level, median [IQR] and worst case per level, the three worst scans per engine with their missed and added regions, and local error maps. & S3 \\ \midrule
\textbf{I12} Sensitivity of decimation and closure across settings and structures. & Closure by decimation for every organ, source and level (S6), the grid comparison on both organs (H6) and the per-configuration analysis (S12). & S6, S12 \\ \bottomrule
\end{longtable}
\section*{What the revision could not do, and why}
\begin{itemize}
\item Real thin and thick reconstructions of the same acquisition: the MSD provides one reconstruction per scan. A dataset with paired reconstructions and expert labels is needed.
\item Multiple annotators: the MSD provides one expert segmentation per scan.
\item Other centres, engines and branching structures (for example hepatic vessels): both organs come from one centre through the MSD, and only TotalSegmentator was run.
\item Clinical tolerance limits: the tolerances in S8 are illustrative.
\end{itemize}
\end{document}
